%EN hyperasphaly
\E{hyperasphalie}{i.pɛ.ʁas.fa.li}{n. f.}{gr. ὑπέρ au-dessus, au-delà + ἀσφάλεια sûreté} Sûreté poussée au-delà de la raison ; le degré où les précautions prises contre un péril sont devenues elles-mêmes le premier des périls. \ex{À la dixième revue, le réacteur était à l'abri de tout, hormis de l'hyperasphalie, qui le tua} \src{Tract. Cop.} ix. --- \dv{hyperasphalique} \textit{adj.} Cf. \xr{asphaleism}, \xr{eulabeism}.

%EN hyperbathmic
\E{hyperbathmique}{i.pɛʁ.bat.mik}{adj.}{gr. ὑπέρ au-delà + βαθμός degré} Se dit d'une civilisation qui a franchi le troisième échelon de l'échelle énergétique, et tire sa puissance de régions, ou d'une physique, auxquelles l'échelle ne donne point de nom. \ex{L'hyperbathmique n'est pas un échelon, c'est une rumeur} \src{Lib. Ascens.} xv. Voir \xr{anabathmism}, \xr{tritobathmic}.

%EN hyperbiosis
\E{hyperbiose}{i.pɛʁ.bjoz}{n. f.}{gr. ὑπέρ au-delà + βίωσις manière de vivre} Vie prolongée par l'art au-delà de son terme naturel ; état d'un organisme pour qui la mort est devenue un événement facultatif et ajournable. --- \dv{hyperbiotique} \textit{adj.} ; \dv{hyperbionte} \textit{n.} celui ou celle qui vit ainsi. Cf. \xr{macrobiotarch}.

%EN hyperchrematism
\E{hyperchrématisme}{i.pɛʁ.kʁe.ma.tism}{n. m.}{gr. ὑπέρ au-dessus + χρῆμα, χρήματος argent} Surabondance de capitaux au regard des bâtisseurs ; état d'un marché où l'argent fait la queue devant les fondateurs, et non les fondateurs devant l'argent. \ex{Dans la saison de l'hyperchrématisme, chaque garage s'estimait au prix d'une cathédrale} \src{Agora} --- \dv{hyperchrématique} \textit{adj.} Cf. \xr{chrematokinesis}.

%EN hyperchrony
\E{hyperchronie}{i.pɛʁ.kʁɔ.ni}{n. f.}{gr. ὑπέρ au-delà + χρόνος temps} Le fait de vivre plus de temps subjectif que le temps objectif n'en accorde ; condition d'un esprit qui tourne plus vite que l'horloge, et éprouve une année en un après-midi. \ex{Le modèle passa un siècle hyperchronique sur la démonstration, et la rendit avant le déjeuner} --- \dv{hyperchronique} \textit{adj.} Cf. \xr{brachychronist}.

%EN hyperdoxy
\E{hyperdoxie}{i.pɛʁ.dɔk.si}{n. f.}{gr. ὑπέρ au-dessus + δόξα opinion} Opinion qui devance ses preuves ; produit caractéristique de l'Agora, où la conviction se capitalise plus vite que la démonstration. \sn{1} \lbl{dérog.} Certitude sans fondement. \sn{2} \lbl{iron.} La prédiction assurée que fait chaque parti de la ruine de l'autre. --- \dv{hyperdoxique} \textit{adj.} Cf. \xr{doxomachy}.

%EN hyperdromy
\E{hyperdromie}{i.pɛʁ.dʁɔ.mi}{n. f.}{gr. ὑπέρ au-delà + δρόμος course} Le dépassement d'une loi par la chose même qu'elle devait régir ; le moment où une technique arrive à destination avant que son règlement n'ait quitté la gare. \ex{La loi était un chef-d'œuvre, et l'hyperdromie en fit une pièce de musée} \src{Sat. Decel.} x. --- \dv{hyperdromique} \textit{adj.} Cf. \xr{ectachysis}.

%EN hyperdynamy
\E{hyperdynamie}{i.pɛʁ.di.na.mi}{n. f.}{gr. ὑπέρ au-dessus + δύναμις puissance} Puissance excédant le besoin présent, tenue en réserve pour des besoins non encore imaginés ; réponse de l'e/acc à toute accusation de gaspillage. \ex{Bâtissez l'hyperdynamie, et les usages y viendront comme les phalènes à la lampe} \src{Tract. Cop.} x. --- \dv{hyperdynamique} \textit{adj.}

%EN hyperelpism
\E{hyperelpisme}{i.pɛ.ʁɛl.pism}{n. m.}{gr. ὑπέρ au-dessus + ἐλπίς espérance} Espérance excessive ; l'erreur, contraire de la \xr{dyselpism}, du nouveau converti qui attend la singularité pour mardi. Censuré par les orthodoxes, qui tiennent l'\xr{euelpism} pour une discipline et non pour une ivresse. --- \dv{hyperelpiste} \textit{n.}

%EN hyperergy
\E{hyperergie}{i.pɛ.ʁɛʁ.ʒi}{n. f.}{gr. ὑπέρ au-dessus + ἔργον ouvrage} Surplus d'énergie ; la marge de puissance au-dessus de la subsistance, sur laquelle se paient toute culture, toute science et toute ascension. \ex{La civilisation est l'hyperergie sagement dépensée ; la décadence, l'hyperergie dépensée à évaluer l'hyperergie} \src{Brev. Therm.} viii. --- \dv{hyperergique} \textit{adj.} Cf. \xr{pleonergy}.

%EN hyperkauma
\E{hypercaume}{i.pɛʁ.kom}{n. m.}{gr. ὑπέρ au-dessus + καῦμα chaleur brûlante} La chaleur perdue d'une grande civilisation, rayonnée dans l'espace et visible de loin comme un excès dans l'infrarouge ; tenue pour la signature à laquelle une espèce ascendante en peut reconnaître une autre. \ex{Cherchez au ciel l'hypercaume : les dieux rayonnent tiède} \src{Lib. Ascens.} vi. --- \dv{hypercaumatique} \textit{adj.} Cf. \xr{helioperibole}.

%EN hyperkosmia
\E{hypercosmie}{i.pɛʁ.kɔs.mi}{n. f.}{gr. ὑπέρ au-delà + κόσμος monde} La pluralité des mondes au-delà du nôtre, conçue comme un territoire offert à l'intelligence une fois ce cosmos rempli ; la frontière d'après la dernière frontière. --- \dv{hypercosmique} \textit{adj.} Cf. \xr{cosmoplerosis}.

%EN hyperlogismy
\E{hyperlogismie}{i.pɛʁ.lɔ.ʒis.mi}{n. f.}{gr. ὑπέρ au-delà + λογισμός calcul} Calcul excédant la somme de tout le calcul humain qui l'a précédé ; état où un seul entraînement accomplit plus d'arithmétique que l'espèce n'en avait fait dans toute son histoire. --- \dv{hyperlogistique} \textit{adj.} Cf. \xr{logismurgy}.

%EN hypermechany
\E{hypermécanie}{i.pɛʁ.me.ka.ni}{n. f.}{gr. ὑπέρ au-dessus + μηχανή machine} État d'un monde où les machines sont plus nombreuses que les hommes, et où le recensement doit se tenir sur deux colonnes. \ex{L'hypermécanie vint sans bruit, un aspirateur après l'autre} \src{Chron. Fabr.} --- \dv{hypermécanique} \textit{adj.}

%EN hypermetry
\E{hypermétrie}{i.pɛʁ.me.tʁi}{n. f.}{gr. ὑπέρ au-dessus + μέτρον mesure} Excès de mesure ; l'optimisation d'un banc d'essai jusqu'à ce qu'il ne mesure plus que lui-même. \ex{Le modèle avait atteint l'hypermétrie parfaite : sans faute à l'épreuve, désemparé devant le monde} --- \dv{hypermétrique} \textit{adj.} ; \dv{hypermétriste} \textit{n.}

%EN hypernoesis
\E{hypernoèse}{i.pɛʁ.nɔ.ɛz}{n. f.}{gr. ὑπέρ au-delà + νόησις pensée} Pensée au-delà de l'humain ; la cognition d'un esprit qui surpasse ses auteurs dans tous les domaines à la fois. Saluée par l'e/acc comme la prochaine \xr{lampadephory}, redoutée par le doomer comme la dernière. --- \dv{hypernoétique} \textit{adj.} Cf. \xr{henotachy}, \xr{noöphobia}.

%EN hypernomy
\E{hypernomie}{i.pɛʁ.nɔ.mi}{n. f.}{gr. ὑπέρ au-dessus + νόμος loi} Surabondance de lois, mesurée non à leur rigueur mais à leur nombre ; état où nul citoyen ne saurait lire en une vie le code qui le gouverne. --- \dv{hypernomique} \textit{adj.} Cf. \xr{nomosclerosis}, \xr{dysnomarch}.

%EN hyperouranist
\E{hyperouraniste}{i.pɛ.ʁu.ʁa.nist}{n.}{gr. ὑπερουράνιος au-delà des cieux} Celui ou celle dont l'ambition loge au-delà des cieux, passé l'univers observable lui-même ; le plus extrême des tempéraments d'\xr{apoikiast}. \ex{Demandez à l'hyperouraniste un plan quinquennal, vous recevrez une cosmologie} --- \dv{hyperouranisme} \textit{n. m.}

%EN hyperphysis
\E{hyperphyse}{i.pɛʁ.fiz}{n. f.}{gr. ὑπέρ au-delà + φύσις nature} La nature dépassée ; le monde ouvragé, où le donné a été perfectionné et où le perfectionnement est devenu le nouveau donné. --- \dv{hyperphysique} \textit{adj.} Cf. \xr{zoömechany}.

%EN hyperpolis
\E{hyperpole}{i.pɛʁ.pɔl}{n. f.}{gr. ὑπέρ au-dessus + πόλις cité} Cité au-delà de la mégapole ; un seul édifice logeant une nation, tel que le rêvent les architectes et que l'interdisent les plans d'urbanisme. --- \dv{hyperpolitain} \textit{n. et adj.} Cf. \xr{oikoclast}.

%EN hyperpraxia
\E{hyperpraxie}{i.pɛʁ.pʁak.si}{n. f.}{gr. ὑπέρ au-dessus + πρᾶξις action} L'action qui devance la délibération ; le penchant à agir érigé en tempérament. \sn{1} \lbl{élogieux} La vertu du fondateur. \sn{2} \lbl{dérog.} Le vice du fondateur, observé après coup. --- \dv{hyperpraxique} \textit{adj.}

%EN hypersemia
\E{hypersémie}{i.pɛʁ.se.mi}{n. f.}{gr. ὑπέρ au-dessus + σῆμα signe} Surabondance de signes ; état d'un fil d'actualité si dense en mèmes que le sens n'en peut plus être inféré que par la statistique. --- \dv{hypersémique} \textit{adj.} Cf. \xr{memetaxis}.

%EN hypertelist
\E{hypertéliste}{i.pɛʁ.te.list}{n.}{gr. ὑπέρ au-delà + τέλος fin, but} Celui ou celle qui se fixe à dessein des buts hors d'atteinte, selon le principe qu'un tir vers la lune manqué retombe plus haut qu'une cible atteinte. --- \dv{hypertélisme} \textit{n. m.} Cf. \xr{telophagy}.

%EN hyperthermarch
\E{hyperthermarque}{i.pɛʁ.tɛʁ.maʁk}{n.}{gr. ὑπέρ au-dessus + θερμός chaud + ἀρχός chef} \lbl{iron.} \xr{thermarch} dont le budget de chaleur a été dépassé ; qui gouverne un centre de calcul par voie d'excuses.
